\scanpage{50}
que \deleted{la famille soit}\inserted{au voisinage} de chaque pt de $\Sigma^{<}(d)\simeq\Sigma^{<}(n)$ $\Sigma^{<}(d)$ soit l'intersection \inserted{dans $\Sigma^{<}(i_n)\simeq\Sigma(0)$} des $\Sigma^{<}(\{i_\alpha,i_n\})_\alpha$ ($\simeq\Sigma^{<}(1)$), et que cette famille soit transversale.

Cela implique la chose analogue pour tout morphisme d'immersion
\[
\Sigma^{<}(d')\longrightarrow\Sigma^{<}(d)
\]
\margin{$d\subset d'$, $d,d'$ cofinaux (ont \unsure{le même} plus grand élément)}
\deleted{\unsure{dans} la structure,}
on considère les $d''$ entre $d$ et $d'$ (donc cofinaux avec $d$) qui sont de la forme $d\cup\{i\}$ ($i\in d'\smallsetminus d$) on regarde le diagramme
\[
(34)
\]
\diagram{\begin{tikzcd}[row sep=large,column sep=large]
&\Sigma^{<}(d\unsure{'}\cup j_0)\arrow[dr,"\text{bordant}"]&\\
\Sigma^{<}(d')\arrow[ur]\arrow[ddr,no head]&&\Sigma^{\unsure{*|<}}(d)\\
&\vdots&\\
&\Sigma^{<}(d\cup j_p)\arrow[uur,no head,"\text{bordant}"']&
\end{tikzcd}}
qui généralise (33), et on \deleted{suppose}\inserted{voit} qu'il satisfait la même condition de transversalité. Car on identifie tous les espaces envisagés, au voisinage des pts en jeu, à des sous-espaces de $\Sigma^{<}(\{\unsure{\alpha}\})\simeq\Sigma^{\unsure{*|<}}(0)$, décrits \inserted{(grâce à b)} comme des intersections partielles des certains sous-espaces bordants de $\Sigma^{<}(\{\unsure{\alpha}\})$, se coupant transversalement\ldots
\mspage{49}
